% !TeX root = ../main.tex

\chapter{Conclusion and Future Work}
\label{ch:conclusion}

\section{Summary of Contributions}
\label{sec:concl-summary}

This thesis asked whether a single-run, teacher-free \emph{average-velocity} residual can replace the EDM diffusion residual in a regional, convection-permitting StormCast forecaster over the Taiwan domain --- delivering the speed a curved diffusion ODE cannot, without sacrificing the sharpness and probabilistic calibration that motivated generative forecasting in the first place --- and how few network evaluations such a residual ultimately needs. Starting from the StormCast two-stage decomposition (a frozen deterministic regression mean plus a generative residual head) trained on the Taiwan RWRF dataset, we adapted \textbf{MeanFlow} --- which learns the \emph{average} velocity over a whole time interval and transports the state across it in one or two network evaluations, with no integration at all --- as the generative residual head. As baselines we retained the EDM diffusion residual it sets out to accelerate and \textbf{FlowCast}, a Conditional-Flow-Matching head that learns the straight-line instantaneous velocity field and samples it with a handful of explicit Euler steps and that MeanFlow generalises exactly. All three were compared head-to-head on the full 2022 hourly validation year.

The concrete contributions were enumerated in Section~\ref{sec:intro-contributions}; the headline finding is that, at a matched $\approx 2\,\text{M}$-sample budget, MeanFlow reaches the diffusion baseline's deterministic and categorical skill class at $K \in \{1, 2\}$ evaluations --- roughly a tenth of the diffusion sampler's cost --- posting the best CRPS (best in the study, though within the single-snapshot noise floor), the best categorical precipitation hit rates (CSI/HSS) at the headline neighbourhood, and the best temperature accuracy of any head, conceding only a spread-driven ensemble-calibration gap at $K = 1$ that the second evaluation closes. The FlowCast baseline meanwhile matches the diffusion CRPS, beats it on wind and precipitation RMSE and on false-alarm ratio at the wider neighbourhoods, and runs $\approx 3.7\times$ faster at 10 Euler steps --- the intermediate point on the same trajectory axis.

\section{Limitations}
\label{sec:concl-limitations}

Several limitations --- together with the failure modes observed during development --- should be kept in mind when interpreting these results.

\paragraph{Single training snapshot.} The headline comparison is taken at a single $\approx 2\,\text{M}$-sample checkpoint per head, where single-snapshot validation noise is $10$--$30\,\%$ (MeanFlow's checkpoint sits essentially on the anchor; Section~\ref{sec:exp-protocol}). The canonical FlowCast run continues to $\approx 13.4\,\text{M}$ samples, and the diffusion baseline likewise has more training available; a converged-budget comparison is not yet in hand.

\paragraph{Single regression mean and single domain.} All three heads are anchored to one frozen regression mean trained on one regional domain (Taiwan RWRF). We cannot assess transfer to other convection-permitting domains (HRRR/CONUS StormCast) without retraining, nor disentangle how much of the flow heads' behaviour depends on the specific quality of this regression mean.

\paragraph{Heavy-rain detection.} The flow heads' conservative, low-false-alarm behaviour costs pixel-exact hits at the heavier rainfall thresholds. The neighbourhood FSS recovers part of this, suggesting a displacement rather than a pure miss, but the deterministic heavy-rain CSI is a genuine weakness at the current budget --- and at the heaviest threshold, five-member ensemble-mean scoring degenerates for every head (averaging suppresses extremes), so the experiment cannot yet rank heads there at all.

\paragraph{Limited probabilistic evaluation.} CRPS is computed from five-member per-hour ensembles --- a small, operationally realistic ensemble in the StormCast tradition --- but a large-ensemble calibration study (rank histograms, spread--skill at 50--100 members), which a cheap few-step sampler makes affordable, has not been run.

\paragraph{Under-dispersion at low NFE (MeanFlow).} At $K = 1$ MeanFlow's deterministic skill is intact but its ensemble spread narrows --- a single average-velocity jump lands each noise draw nearly on the conditional mean --- so its aggregate CRPS lags its own $K = 2$ setting, most visibly on the wind channels, and the deficit \emph{compounds under autoregression}: over the six-hour rollout the $K = 1$ CRPS ($0.5501$) trails $K = 2$ ($0.5207$) by $\approx 5.6\,\%$; at $K = 2$ MeanFlow in fact attains the best aggregate CRPS of all heads ($0.5207$ against the ODE heads' $0.5518$/$0.5537$), so no residual gap against them remains. The second evaluation restores most of the spread, which is why $K = 2$ is the recommended operating point; a per-member rank-histogram diagnosis has not yet been run.

\paragraph{No distilled-diffusion comparator.} We compare the flow heads against the full $35$-NFE diffusion sampler, not against a distilled few-step diffusion student. MeanFlow demonstrates that the $1$--$2$-NFE regime is reachable in a single run with no teacher; whether distillation could carry the diffusion baseline's heavy-rain CSI into that same regime is an open question this thesis does not answer.

\paragraph{FlowCast \texttt{t2m} rollout drift.} FlowCast's temperature error grows fastest of the cleaned heads under rollout, reaching \SI{1.73}{\kelvin} at $+6$\,h against the diffusion baseline's \SI{1.60}{\kelvin} (Figure~\ref{fig:rollout}) --- the autoregressive signature of its weaker single-step \texttt{t2m}. Notably MeanFlow does \emph{not} inherit this drift despite sharing FlowCast's backbone and loss weights (its \texttt{t2m} stays below the diffusion baseline at \SI{1.57}{\kelvin}), so the drift is a property of the learned trajectory rather than the architecture; a per-channel weight that does not starve \texttt{t2m}, or a short teacher-forcing warmup, is the natural fix.

\paragraph{EMA evaluation pitfall.} Loading the online \texttt{checkpoint\_*.pt} training-state dump instead of the EMA shadow (\texttt{ema\_state.pt}) produces visibly noisier flow-head samples. This is an evaluation pitfall rather than a model failure: the head-to-head harness loads the online \texttt{.mdlus} checkpoints for FlowCast and MeanFlow \emph{symmetrically}, so the A/B comparison is unaffected, and the EMA-versus-online gap shrinks as training proceeds.

\section{Future Work}
\label{sec:concl-future}

The most promising directions, several of which are scoped concretely in the accompanying project to-do, are:
\begin{enumerate}
  \item \textbf{Train to convergence and re-score.} Run all three heads to a matched, converged sample budget and repeat the full metric suite, smoothing over several validation snapshots to beat down the $10$--$30\,\%$ noise. This is the single most important step to turn the provisional headline into a defensible claim.
  \item \textbf{Isolate the loss recipe.} Re-train the EDM diffusion baseline with the same rain-aware channel weight $w_{\text{pr}}$ the flow heads use, splitting the recipe's contribution from the generative objective's. The speed and NFE results are untouched by this experiment; the skill gaps are its target.
  \item \textbf{Close the heavy-rain gap.} Test an explicit tail-quantile or focal loss on \texttt{qpepre}, a higher $w_{\text{pr}}$, and a stronger \texttt{asinh}/\texttt{log1p} dynamic-range transform, targeting the $5$--$10$\,mm\,h$^{-1}$ CSI deficit without inflating false alarms.
  \item \textbf{Restore low-NFE rollout spread.} MeanFlow under-disperses at $K = 1$ and the deficit compounds over autoregressive rollouts (a residual CRPS gap survives even at $K = 2$); candidate fixes --- a stochastic single-segment sampler, temperature on the noise draw, per-step noise re-injection during rollout, or a small spread-aware loss term --- would make the true 1--2-NFE regime probabilistically competitive end to end, diagnosed with per-member rank histograms.
  \item \textbf{Large-ensemble calibration.} Exploit the flow heads' low per-sample cost to run 50--100-member ensembles and report CRPS, rank histograms, and spread--skill --- the uncertainty-quantification payoff that cheap sampling unlocks. At MeanFlow's two evaluations, a 100-member 24-hour cycle costs $4{,}800$ network evaluations --- about a ninth of what the 35-NFE diffusion sampler needs for a 50-member one --- so ensemble size, not the sampler, becomes the budget knob.
  \item \textbf{Transfer to other domains.} Apply the same single-run flow recipes (FlowCast and MeanFlow) to another convection-permitting regression mean (e.g.\ an HRRR-trained StormCast) to measure how much of the pipeline is Taiwan-specific.
\end{enumerate}

\section{Closing Remarks}
\label{sec:concl-closing}

Generative modelling is almost certainly the right inductive bias for short-range convective-scale forecasting; the obstacle to deploying it has been the multi-step diffusion sampler. This thesis shows that, on a real regional dataset, the path a generative model takes from noise to weather is a dial worth turning twice --- first straightened by a single-run straight-line flow residual (most of the diffusion skill at roughly a quarter of the cost, with no teacher), then jumped by an average-velocity residual that compresses the sampler to one or two evaluations, a tenth of the diffusion cost, at the study's best CRPS and categorical precipitation hit rates and trading only a modest one-evaluation spread gap. The remaining gaps --- heavy-rain hits, FlowCast's \texttt{t2m} drift, one-evaluation spread, convergence at a larger budget --- are concrete and addressable. Making sharp, probabilistic, regional convective forecasting fast enough to run between radar volumes is within reach, and single-run flow residuals --- straight-line or average-velocity --- are a credible path to it.
